\documentclass[UTF8,12pt]{ctexart}

\usepackage[a4paper,margin=2.5cm]{geometry}
\usepackage{amsmath,amssymb}
\usepackage{booktabs}
\usepackage{array}
\usepackage{graphicx}
\usepackage{float}
\usepackage{caption}
\usepackage{setspace}
\usepackage[colorlinks=true,linkcolor=black,citecolor=black,urlcolor=black]{hyperref}

\captionsetup{font=small}
\setlength{\tabcolsep}{4pt}
\renewcommand{\arraystretch}{1.12}
\onehalfspacing

\title{绿色信贷、高煤炭地区转型与地方承接能力：\\基于连续煤炭暴露度和晋蒙路径差异的证据}
\author{}
\date{\today}

\begin{document}

\maketitle

\begin{abstract}
本文考察绿色信贷是否推动高煤炭地区低碳转型，并进一步识别其作用到哪一层。研究以 2012 年《绿色信贷指引》为全国性政策冲击，结合政策前煤炭暴露度，构造连续暴露度双重差分。结果显示，绿色信贷确实对高煤炭地区产生影响，但其首先表现为能源使用强度下降、煤炭使用强度下降和污染治理改善。绿色全要素生产率和全局曼奎斯特--卢恩伯格指数未显著变化，说明当前证据不足以支持绿色生产率跃升的强结论。煤炭消费占比显著下降，但二氧化碳排放、火电装机占比和火电发电占比没有同步稳定调整，表明煤炭占比变化仍属于约束性治理向结构性治理的过渡性调整。进一步的三重差分结果显示，低碳禀赋本身与后期风光增长存在一定关系，但并未自动放大绿色信贷约束向新能源替代的转化。山西和内蒙古的政策文本证据展示了两类地方承接路径：山西更偏向存量煤炭体系清洁化，内蒙古更容易形成新能源增量承接叙事。本文将绿色信贷理解为高煤炭地区转型的外部约束和治理触发机制；结构替代还需要额外的地方承接条件。
\end{abstract}

\noindent\textbf{关键词：}绿色信贷；煤炭暴露；污染治理；能源结构；政策文本；晋蒙转型

\newpage

\section{概念框架、数据与实证设计}

\subsection{研究问题与核心概念}

本文聚焦绿色信贷在高煤炭地区的作用边界。核心命题是：绿色信贷首先约束高煤炭地区，结构性改造能力需要进一步检验。它作为金融治理工具，通过信贷准入、融资约束、环境风险审查和项目合规要求，对既有高耗能、高排放活动施加外部约束；这种约束可以先触发治理性调整，但结构性低碳替代是否随后发生，取决于地方低碳禀赋、项目建设、电网消纳、外送通道和政策承接能力。

本文将绿色信贷作用区分为约束性治理和结构性治理。约束性治理指绿色信贷通过信贷准入、融资约束、环境风险审查和项目合规要求，对高煤炭地区既有高耗能、高排放活动形成外部约束，从而推动能源使用强度下降、煤炭使用强度下降和污染排放下降。它主要作用于既有高煤炭体系内部的运行方式和行为边际，不要求能源资本结构立即发生替代。

结构性治理指绿色信贷或转型金融进一步改变高煤炭地区的能源投资方向、发电结构、碳排放路径和低碳项目体系，使地方发展从煤炭依赖和火电锁定转向新能源替代、低碳基础设施和持续的政策承接。它同时改变煤炭使用方式和煤炭体系在地方发展中的主轴地位。

二者的核心区别在于：约束性治理改变的是运行强度，结构性治理改变的是结构配置；约束性治理降低的是高煤炭体系的运行强度，结构性治理改变的是高煤炭体系的再生产方式。约束性治理可以先发生，但结构性治理不一定自动随后发生。本文聚焦绿色信贷有效到哪一层，以及为什么无法自动进入下一层。

\subsection{研究对象与数据来源}

研究起点是 2012 年《绿色信贷指引》形成的全国性政策冲击。样本为中国省级年度面板，主体样本覆盖 2000--2023 年。基础宏观变量来自《中国统计年鉴》、Wind、CEIC 和 EPS 等常用省级数据库；能源消费与能源结构变量主要来自《中国能源统计年鉴》的地区能源平衡表；电力装机与发电结构变量来自《中国电力统计年鉴》的分地区装机和发电量表；环境变量来自《中国环境统计年鉴》和生态环境统计年报；二氧化碳排放来自 CEADs 或 EDGAR 省级排放清单。绿色全要素生产率和全局曼奎斯特--卢恩伯格指数沿用既有测算口径或已有数据，并在实证中作为能源强度压降之外的补充结果。

本文进一步构建山西和内蒙古两省政策文本数据。山西政策文本来自山西省人民政府政策文件库，覆盖省委文件、省政府令、省政府文件、省政府办公厅文件和部门规范性文件。内蒙古政策文本来自内蒙古自治区人民政府政府信息公开目录，覆盖政府文件、政府令、行政规范性文件及其财政金融、自然资源能源、工业交通、生态环境等分类。文本抓取后保留标题、发布日期、发文字号、正文、正文长度和政策分类，并汇总为省份--年份层面的政策注意力指标。山西共获得 3874 条政策记录，其中 3871 条成功抓取正文；内蒙古共获得 2974 条政策记录，其中 2972 条成功抓取正文。政策文本变量用于晋蒙路径差异和机制补充，不进入全国连续 DID 的主识别。

\subsection{核心解释变量：政策冲击与煤炭暴露}

本文将 2012 年绿色信贷政策作为全国性政策冲击。记 $\text{政策后}_{t}$ 为年份不早于 2012 的虚拟变量。为避免将山西和内蒙古直接设定为唯一处理组，本文使用政策前煤炭暴露度刻画各省受绿色信贷约束的处理强度。核心煤炭暴露度为 2008--2011 年政策前煤炭终端消费占能源终端消费的平均水平，记为 $\text{煤炭暴露}_{i}$。主解释变量定义为：
\[
    \begin{aligned}
    \text{政策后煤炭暴露}_{it}
    &= \text{政策后}_{t} \times \text{煤炭暴露}_{i}.
    \end{aligned}
\]
该设定把全国绿色信贷政策冲击与省份间政策前煤炭依赖差异结合起来。系数衡量的是：2012 年之后，政策前煤炭暴露度更高的省份是否出现更明显的效率、污染治理和结构变量变化。

在低碳替代部分，本文进一步使用后期政策窗口和低碳禀赋变量构造三重交互项。由于风电、光伏等新能源装机的大规模扩张与 2012 年绿色信贷政策之间存在时间错位，本文将 2016 年以后作为低碳替代承接的后期窗口：
\[
    \begin{aligned}
    \text{后期煤炭暴露}_{it}
    &= \text{后期政策}_{t} \times \text{煤炭暴露}_{i}.
    \end{aligned}
\]
低碳禀赋主要用政策前非火电装机占比、早期风电装机占比和早期风电发电占比表示。地方承接能力是指一个高煤炭地区把绿色信贷形成的外部约束，转化为低碳项目建设、能源系统调整和新能源替代的能力。完整的地方承接能力超出单纯的低碳资源禀赋和政策文本中的绿色表述，由资源基础、电网消纳、项目建设和政策议程共同构成。考虑到目前缺少完整的外送电量、电网容量和弃风弃光长面板数据，本文在全国连续 DID 中采用弱版低碳禀赋指数：
\[
    \begin{aligned}
    \text{低碳禀赋}_{i}
    &= \text{均值}\big[
    z(\text{政策前非火电装机}_{i}), \\
    &\qquad
    z(\text{早期风电装机}_{i}),
    z(\text{早期风电发电}_{i})
    \big].
    \end{aligned}
\]
其中 $\text{政策前非火电装机}_{i}$ 为政策前非火电装机占比，$\text{早期风电装机}_{i}$ 为早期风电装机占比，$\text{早期风电发电}_{i}$ 为早期风电发电占比。该指数刻画先验承接条件或低碳前置条件，不等同于完整的地方承接能力。为避免将政策后的转型结果误作为解释变量，全国模型仅使用政策前或早期形成的低碳禀赋作为调节变量。政策文本变量用于刻画山西和内蒙古在政策后的承接表现，即地方政府是否持续将绿色金融、煤炭清洁治理和新能源替代纳入政策议程。由此，本文区分先验承接条件和后验承接表现：前者用于检验绿色信贷作用的条件性，后者用于展示高煤炭地区转型路径的地方差异。进一步构造如下三重交互：
\[
    \begin{aligned}
    \text{三重交互}_{it}
    &= \text{政策后}_{t} \times \text{煤炭暴露}_{i}
    \times \text{低碳禀赋}_{i}.
    \end{aligned}
\]
该设定用于检验高煤炭暴露地区的低碳前置条件是否会改变绿色信贷约束与新能源替代之间的关系。

\subsection{变量构造与概念对应}

本文按照约束性治理、过渡性调整和结构性治理组织被解释变量。下表给出概念、变量和证据定位之间的对应关系。

\begin{table}[htbp]
    \centering
    \caption*{表 X：约束性治理、过渡性调整与结构性治理的变量对应}
    \label{tab:conceptmap}
    \resizebox{\textwidth}{!}{%
    \begin{tabular}{p{3.1cm}p{4.2cm}p{4.5cm}p{5.1cm}}
        \hline
        概念层次 & 理论含义 & 对应变量 & 证据定位 \\
        \hline
        约束性治理：能源强度压降 & 在既有高煤炭体系内部降低能源和煤炭投入强度 & 五类能源终端消费强度、煤炭终端消费强度 & 主证据。二者显著下降，说明绿色信贷首先压低能源和煤炭使用强度。 \\
        约束性治理：污染治理 & 高污染活动受到融资和环境风险约束，末端治理增强 & 工业二氧化硫、氮氧化物、工业固体废物，附带颗粒物、工业废水 & 主证据。SO$_2$、NOx、固废显著下降，PM 和废水不稳。 \\
        过渡性调整 & 煤炭使用比例下降，但尚未证明能源系统重构 & 五类能源口径煤炭占比、原煤炭消费占比 & 较强证据。煤炭占比下降说明约束性治理向结构性治理出现边际过渡，但还不足以证明完整去煤。 \\
        结构性治理：深层去煤 & 火电资本、发电结构和碳排放同步调整 & 二氧化碳排放对数、火电装机占比、火电发电占比 & 边界证据。目前不显著，说明结构性治理没有自动发生。火电发电占比存在前趋势问题，只作为结构锁定的描述性证据。 \\
        结构性治理：低碳替代 & 新能源发电、装机和低碳项目体系扩张 & 非火电发电占比、风电发电占比、风光发电占比、风电装机占比 & 条件性证据。存在时间错位和承接能力约束，不能被解释为绿色信贷的机械结果。 \\
        地方承接能力与后验承接表现 & 地方是否有资源、项目、基础设施、消纳条件和政策连续性，把金融约束转化为结构替代 & 政策文本中的绿色金融、煤炭清洁治理、污染治理、新能源替代词频和覆盖比例；晋蒙路径；领导更替节点 & 机制和案例证据，不进入全国连续 DID 的主识别。 \\
        \hline
    \end{tabular}
    }
\end{table}

第一类是约束性治理变量，包括五类能源终端消费强度、煤炭终端消费强度、绿色全要素生产率水平和全局曼奎斯特--卢恩伯格指数。其中能源强度变量以能源或煤炭终端消费除以地区生产总值构造，用以刻画高煤炭体系内部的强度压降。绿色全要素生产率和全局曼奎斯特--卢恩伯格指数用于检验这种强度压降是否进一步表现为广义绿色生产率变化。

第二类是污染治理变量，包括工业二氧化硫排放、氮氧化物排放、颗粒物排放、工业固体废物和工业废水排放。这类变量用于检验绿色信贷政策是否首先通过污染治理和末端减排体现出来。

第三类是过渡性调整和深层结构性治理变量。五类能源口径煤炭占比和原煤炭消费占比刻画煤炭使用比例变化，是约束性治理向结构性治理过渡的信号。二氧化碳排放对数、火电装机占比和火电发电占比刻画碳排放路径、火电资本存量和发电结构，用于判断煤炭占比下降是否同步转化为深层去煤。

第四类是低碳替代变量，包括非火电发电占比、风电发电占比、风光发电占比和风电装机占比。该组变量用于检验绿色信贷政策是否在低碳禀赋较强的地区进一步转化为新能源替代。

第五类是政策文本变量。本文将每篇政策文件的标题、主题和正文合并后，对四类关键词进行词频统计：绿色金融、煤炭清洁治理、污染治理和新能源替代。对每一类主题，先在文档层面统计关键词出现次数，再汇总到省份--年份层面。以新能源主题为例：
\[
    \begin{aligned}
    \text{新能源词频}_{it}
    &= \sum_{d \in (i,t)}
    \sum_{k \in K_R} Count(k,d),
    \end{aligned}
\]
其中 $K_R$ 包括 ``新能源''、``风电''、``光伏''、``太阳能''、``可再生能源''、``绿电''、``储能''、``源网荷储''、``外送''、``特高压''、``绿色电力''、``氢能'' 和 ``绿氢'' 等关键词。为避免年份间文件数量和正文长度差异影响解释，本文构造每万字词频：
\[
    \begin{aligned}
    \text{主题每万字词频}_{it}
    &= \frac{\text{主题词数}_{it}}{\text{正文字符数}_{it}}
    \times 10000,
    \end{aligned}
\]
并构造主题政策覆盖比例：
\[
    \begin{aligned}
    \text{主题政策覆盖比例}_{it}
    &= \frac{\#\{d:\text{主题词数}_{idt}>0\}}
    {\text{政策文件数}_{it}}.
    \end{aligned}
\]
这些变量刻画地方政府政策文本注意力，而不直接等同于政策支持强度。

\subsection{模型设定}

主模型采用省份和年份双向固定效应：
\[
    \begin{aligned}
    \text{结果变量}_{it}
    &= \beta \text{政策后煤炭暴露}_{it}
    + \gamma \text{控制变量}_{it}
    + \mu_i + \lambda_t + \varepsilon_{it},
    \end{aligned}
\]
其中 $\text{结果变量}_{it}$ 表示效率、污染、结构或低碳替代变量；$\text{控制变量}_{it}$ 包括地区生产总值、人口、产业结构、城镇化率、环保财政支出占比和市场化指数等控制变量；$\mu_i$ 和 $\lambda_t$ 分别表示省份固定效应和年份固定效应。标准误聚类到省份层面。

连续暴露度事件研究模型为：
\[
    \begin{aligned}
    \text{结果变量}_{it}
    &= \sum_{\tau \neq -1}\beta_{\tau}
    \left(\text{煤炭暴露}_{i} \times 1[t-2012=\tau]\right) \\
    &\quad + \gamma \text{控制变量}_{it}
    + \mu_i + \lambda_t + \varepsilon_{it},
    \end{aligned}
\]
其中 2011 年作为基准期。该模型用于检验政策前趋势，并观察政策后系数的动态变化。

低碳禀赋调节模型为：
\[
    \begin{aligned}
    \text{结果变量}_{it}
    &= \beta \left(\text{政策后}_{t} \times \text{煤炭暴露}_{i}\right) \\
    &\quad + \eta \left(\text{政策后}_{t} \times \text{低碳禀赋}_{i}\right) \\
    &\quad + \theta \left(\text{政策后}_{t} \times \text{煤炭暴露}_{i}
    \times \text{低碳禀赋}_{i}\right) \\
    &\quad + \gamma \text{控制变量}_{it}
    + \mu_i + \lambda_t + \varepsilon_{it}.
    \end{aligned}
\]
其中 $\text{政策后}_{t} \times \text{低碳禀赋}_{i}$ 为低碳禀赋的低阶交互项；$\text{煤炭暴露}_{i} \times \text{低碳禀赋}_{i}$ 由于不随时间变化，被省份固定效应吸收。该模型不被解释为直接的新能源替代因果机制，而用于检验低碳替代是否具有条件性和时间错位。

晋蒙政策文本路径采用两省描述性固定效应模型：
\[
    \begin{aligned}
    \text{政策文本结果}_{it}
    &= \delta (\text{内蒙古}_{i} \times \text{政策后}_{t})
    + \mu_i + \lambda_t + \varepsilon_{it}.
    \end{aligned}
\]
该模型仅用于比较 2012 年后内蒙古相对山西的政策文本注意力变化。由于样本只有两个省，该模型作为案例路径证据，不构成全国意义上的因果识别。

\section{实证结果：从约束性治理到结构性治理边界}

\subsection{绿色信贷与能源强度压降：约束性治理的第一层证据}

第一层检验识别绿色信贷是否触发约束性治理。表 \ref{tab:efficiency} 将能源强度压降作为绿色信贷约束进入高煤炭体系后的直接响应。核心解释变量“政策后煤炭暴露”在五类能源终端消费强度上的系数为 $-0.2895$，并在 5\% 水平显著；在煤炭终端消费强度上的系数为 $-0.4262$，并在 1\% 水平显著。该结果表明，2012 年之后，政策前煤炭暴露度越高的省份，能源投入强度和煤炭终端消费强度下降越明显。

绿色全要素生产率和全局曼奎斯特--卢恩伯格指数上的估计系数分别为 $-1.0498$ 和 $-0.1342$，均未达到常规显著性水平。这一结果要求对效率改善作更精确的界定。绿色信贷冲击在当前数据中主要体现为能源和煤炭使用强度下降，尚不足以支持绿色全要素生产率显著提高的强结论。因此，本文将第一层发现界定为高煤炭体系内部的强度压降；广义生产率跃升的证据不足。

\begin{table}[htbp]
    \centering
    \caption{绿色信贷与能源强度压降：约束性治理的第一层证据}
    \label{tab:efficiency}
    \resizebox{\textwidth}{!}{%
    \begin{tabular}{lcccc}
        \hline
        & 绿色全要素生产率 & 全局ML指数 & 五类能源强度 & 煤炭终端强度 \\
        \hline
        政策后煤炭暴露 & -1.0498 & -0.1342 & -0.2895** & -0.4262*** \\
        & (0.7528) & (0.0830) & (0.1190) & (0.0776) \\
        样本量 & 510 & 510 & 480 & 480 \\
        拟合优度 & 0.7644 & 0.2719 & 0.9592 & 0.9237 \\
        省份固定效应 & 是 & 是 & 是 & 是 \\
        年份固定效应 & 是 & 是 & 是 & 是 \\
        \hline
    \end{tabular}
    }
    \begin{flushleft}
    \footnotesize 注：括号内为按省份聚类的稳健标准误。***、**、* 分别表示在 1\%、5\%、10\% 水平显著。
    \end{flushleft}
\end{table}

\subsection{绿色信贷与污染治理：约束性治理的第二层证据}

第二层检验考察约束性治理是否进一步体现为污染治理。表 \ref{tab:pollution} 报告工业二氧化硫和氮氧化物排放的估计结果。“政策后煤炭暴露”对工业二氧化硫的系数为 $-72.5887$，在 1\% 水平显著；对氮氧化物排放的系数为 $-43.9915$，在 5\% 水平显著。附加污染结果中，工业固体废物的系数为 $-47.2683$，在 5\% 水平显著；颗粒物和工业废水系数为负，但未达到常规显著性水平。

这组结果支持一种有边界的污染治理效应：绿色信贷对不同污染物的影响并不均衡，其更稳定地作用于与高煤炭工业生产和能源燃烧联系更紧密的大气污染物与工业固废。该证据表明，绿色信贷首先通过融资约束和环境风险审查强化既有高煤炭体系内部的污染治理，尚未直接说明能源系统结构已经改变。

\begin{table}[htbp]
    \centering
    \caption{绿色信贷与污染治理：约束性治理的第二层证据}
    \label{tab:pollution}
    \resizebox{\textwidth}{!}{%
    \begin{tabular}{lccccc}
        \hline
        & 工业SO$_2$ & 氮氧化物 & 颗粒物 & 工业固废 & 工业废水 \\
        \hline
        政策后煤炭暴露 & -72.5887*** & -43.9915** & -10.1339 & -47.2683** & -40594.6149 \\
        & (19.6996) & (16.8916) & (28.0627) & (20.0521) & (24203.0698) \\
        样本量 & 510 & 336 & 154 & 510 & 510 \\
        拟合优度 & 0.8812 & 0.9223 & 0.8225 & 0.7837 & 0.9108 \\
        省份固定效应 & 是 & 是 & 是 & 是 & 是 \\
        年份固定效应 & 是 & 是 & 是 & 是 & 是 \\
        \hline
    \end{tabular}
    }
    \begin{flushleft}
    \footnotesize 注：括号内为按省份聚类的稳健标准误。***、**、* 分别表示在 1\%、5\%、10\% 水平显著。
    \end{flushleft}
\end{table}

\subsection{从约束性治理到结构性治理：煤炭占比、碳排放与火电结构}

第三层检验区分约束性治理的外溢和结构性治理的完成。表 \ref{tab:structure} 将煤炭占比、碳排放和火电结构同时纳入结果变量。“政策后煤炭暴露”对五类能源口径煤炭占比的系数为 $-0.2495$，在 1\% 水平显著；对原煤炭消费占比的系数为 $-0.1900$，同样在 1\% 水平显著。这说明绿色信贷政策之后，高煤炭暴露地区的煤炭使用占比确实出现下降。

煤炭消费占比的下降说明绿色信贷的约束效应已经外溢到能源使用结构，但这一变化仍停留在过渡性调整层面。更深层的结构变量提供了边界条件：“政策后煤炭暴露”对二氧化碳排放的系数为 $-0.0982$，未显著；对火电装机占比的系数为 $0.0517$，未显著；对火电发电占比的系数为 $0.0888$，也未显著。由于碳排放、火电装机占比和火电发电占比没有同步稳定变化，本文不能将煤炭占比下降直接解释为完整的结构性去煤。换言之，绿色信贷可以改变高煤炭体系的运行强度和部分能源消费比例，但尚未显示出自动重构火电资本存量和发电结构的能力。

\begin{table}[htbp]
    \centering
    \caption{从约束性治理到结构性治理：煤炭占比、碳排放与火电结构}
    \label{tab:structure}
    \resizebox{\textwidth}{!}{%
    \begin{tabular}{lccccc}
        \hline
        & 煤炭占比 & 煤炭占比 & 二氧化碳排放 & 火电装机 & 火电发电 \\
        & 五类能源口径 & 原口径 & 对数 & 占比 & 占比 \\
        \hline
        政策后煤炭暴露 & -0.2495*** & -0.1900*** & -0.0982 & 0.0517 & 0.0888 \\
        & (0.0754) & (0.0651) & (0.1034) & (0.0945) & (0.0754) \\
        样本量 & 480 & 480 & 510 & 510 & 507 \\
        拟合优度 & 0.9224 & 0.9382 & 0.9904 & 0.9368 & 0.9478 \\
        省份固定效应 & 是 & 是 & 是 & 是 & 是 \\
        年份固定效应 & 是 & 是 & 是 & 是 & 是 \\
        \hline
    \end{tabular}
    }
    \begin{flushleft}
    \footnotesize 注：括号内为按省份聚类的稳健标准误。***、**、* 分别表示在 1\%、5\%、10\% 水平显著。
    \end{flushleft}
\end{table}

连续暴露度事件研究显示，不同结果变量的政策前趋势并不完全一致。煤炭消费占比的政策前系数围绕零波动，政策后逐步转负，长期项为 $-0.3685$，在 1\% 水平显著，提供了较强的动态支持。绿色全要素生产率和全局曼奎斯特--卢恩伯格指数的政策前趋势未显示明显系统性差异，但政策后系数缺乏稳定显著性，说明绿色信贷并未带来即时的绿色生产率跃升。工业二氧化硫在政策后持续下降，事件时间 0、1、2、3 和 4 期后的系数分别为 $-15.5398$、$-23.2798$、$-32.7401$、$-45.8361$ 和 $-126.1353$，且均显著；但政策前系数已有一定下行迹象，因此其动态结果应被理解为对污染治理效应的辅助支持，不能作为唯一识别依据。火电发电占比存在明显政策前动态差异，风光发电占比又受数据起点限制，不能作为严格事件研究证据。由此，本文将能源强度压降、污染治理和煤炭占比调整作为主要政策响应，将火电结构和新能源替代作为转型边界和路径差异的补充证据。

\subsection{结构性治理的条件：低碳禀赋、时间错位与新能源替代}

第四层检验解释结构性治理为何没有自动发生。表 \ref{tab:lowcarbon} 使用弱版低碳禀赋指数进行检验，该指数由政策前非火电装机占比、早期风电装机占比和早期风电发电占比标准化后取均值得到。模型同时加入“后期煤炭暴露”、“后期低碳禀赋”和“三重交互”。“后期低碳禀赋”解释低碳禀赋较强的省份在 2016 年后是否普遍出现风光或非火电增长；“三重交互”则解释在高煤炭暴露省份中，低碳禀赋是否进一步放大绿色信贷约束向新能源替代的转化。

结果显示，低碳禀赋本身与后期风光增长存在一定关系：“后期低碳禀赋”对风光发电占比的系数为 $0.0908$，在 10\% 水平显著。但三重交互项并未支持绿色信贷约束在低碳禀赋地区自动转化为新能源替代的强机制。相反，“三重交互”在风电发电占比和风光发电占比上的系数分别为 $-0.0927$ 和 $-0.2363$，均在 5\% 水平显著。该结果说明早期低碳基础并不等于结构性治理能力。新能源替代还受到项目建设周期、电网消纳、外送通道、资源开发边际和地方政策连续性的共同制约。

低碳禀赋只是结构性治理的前置条件，尚不足以构成结构性治理的充分条件。真正的地方承接能力同时包含风光基础、项目建设、并网消纳、外送通道和持续的低碳项目体系，核心在于地方能否将金融约束转化为能源结构重组。

\begin{table}[htbp]
    \centering
    \caption{结构性治理的条件：低碳禀赋、时间错位与新能源替代}
    \label{tab:lowcarbon}
    \resizebox{\textwidth}{!}{%
    \begin{tabular}{lcccc}
        \hline
        & 非火电发电占比 & 风电发电占比 & 风光发电占比 & 风电装机占比 \\
        \hline
        后期煤炭暴露 & -0.2231* & 0.0024 & -0.0919 & 0.0273 \\
        & (0.1164) & (0.0240) & (0.0724) & (0.0675) \\
        后期低碳禀赋 & 0.1072 & 0.0309 & 0.0908* & 0.0246 \\
        & (0.0695) & (0.0182) & (0.0461) & (0.0329) \\
        三重交互 & -0.1757 & -0.0927** & -0.2363** & -0.1077 \\
        & (0.1514) & (0.0402) & (0.1093) & (0.0821) \\
        样本量 & 507 & 330 & 330 & 288 \\
        拟合优度 & 0.9514 & 0.8916 & 0.8921 & 0.9311 \\
        省份固定效应 & 是 & 是 & 是 & 是 \\
        年份固定效应 & 是 & 是 & 是 & 是 \\
        \hline
    \end{tabular}
    }
    \begin{flushleft}
    \footnotesize 注：“三重交互”为“后期政策 $\times$ 煤炭暴露 $\times$ 低碳禀赋”。模型同时控制“后期煤炭暴露”和“后期低碳禀赋”；“煤炭暴露 $\times$ 低碳禀赋”被省份固定效应吸收。低碳禀赋由政策前非火电装机占比、早期风电装机占比和早期风电发电占比标准化后取均值得到。括号内为按省份聚类的稳健标准误。***、**、* 分别表示在 1\%、5\%、10\% 水平显著。
    \end{flushleft}
\end{table}

\subsection{地方承接能力与治理路径分化：山西和内蒙古的政策文本证据}

山西和内蒙古用于展示同样高煤炭暴露下的地方承接路径分化，不承担全国主识别任务。本文将两省核心结果变量按政策前、2012--2016 年和 2017--2023 年三个阶段求均值，并与政策文本强度结合。阶段均值显示，两省煤炭占比均有所下降。山西煤炭占比从政策前的 $0.6431$ 下降到后期的 $0.4305$，内蒙古从 $0.6097$ 下降到 $0.4451$。火电发电占比仍然保持较高水平，山西后期均值为 $0.8511$，内蒙古为 $0.8110$。风光发电占比在后期上升，山西为 $0.1287$，内蒙古为 $0.1738$，内蒙古的增量替代特征更明显。

政策文本证据进一步揭示地方政策注意力的差异。山西的煤炭清洁治理词频长期处于较高水平，政策前每万字词频为 $17.7787$，2012--2016 年为 $10.6072$，2017--2023 年为 $5.5123$。这一趋势显示山西政策叙事长期围绕煤炭安全、先进产能、煤电改造和污染治理展开。内蒙古的新能源词频在后期明显抬升，从政策前的 $1.0490$ 增至 2012--2016 年的 $1.7436$，并在 2017--2023 年达到 $6.2122$。这一变化与内蒙古风光资源、外送通道和新能源基地建设条件相匹配。

两省文本路径回归提供了后验承接表现证据。2012 年之后，内蒙古相对山西的绿色金融每万字词频提高 $0.3062$，在 10\% 水平显著；绿色金融政策覆盖比例提高 $0.0382$，在 5\% 水平显著。煤炭清洁治理每万字词频提高 $9.3511$，政策覆盖比例提高 $0.1361$，二者均在 1\% 水平显著。新能源文本强度系数为 $1.2466$，新能源政策覆盖比例系数为 $0.0102$，方向为正但未达到常规显著性水平。该结果支持一种案例层面的解释：内蒙古在 2012 年后相对山西更明显强化绿色金融和煤炭清洁治理表述，新能源叙事也有上升方向，但其统计稳定性弱于煤炭清洁治理。

\begin{table}[htbp]
    \centering
    \caption{地方承接能力与治理路径分化：晋蒙政策文本证据}
    \label{tab:textpath}
    \resizebox{\textwidth}{!}{%
    \begin{tabular}{lcccc}
        \hline
        & 绿色金融 & 煤炭清洁治理 & 污染治理 & 新能源替代 \\
        & 每万字词频 & 每万字词频 & 每万字词频 & 每万字词频 \\
        \hline
        内蒙古 $\times$ 政策后 & 0.3062* & 9.3511*** & 0.1263 & 1.2466 \\
        & (0.1577) & (2.4313) & (0.7338) & (1.6355) \\
        样本量 & 44 & 44 & 44 & 44 \\
        拟合优度 & 0.6552 & 0.8905 & 0.6290 & 0.6534 \\
        省份固定效应 & 是 & 是 & 是 & 是 \\
        年份固定效应 & 是 & 是 & 是 & 是 \\
        \hline
        & 绿色金融 & 煤炭清洁治理 & 污染治理 & 新能源替代 \\
        & 政策覆盖比例 & 政策覆盖比例 & 政策覆盖比例 & 政策覆盖比例 \\
        \hline
        内蒙古 $\times$ 政策后 & 0.0382** & 0.1361*** & 0.0299 & 0.0102 \\
        & (0.0145) & (0.0398) & (0.0203) & (0.0167) \\
        样本量 & 44 & 44 & 44 & 44 \\
        拟合优度 & 0.8581 & 0.8585 & 0.8402 & 0.8879 \\
        省份固定效应 & 是 & 是 & 是 & 是 \\
        年份固定效应 & 是 & 是 & 是 & 是 \\
        \hline
    \end{tabular}
    }
    \begin{flushleft}
    \footnotesize 注：样本限于山西和内蒙古，标准误为稳健标准误。该表用于描述两省政策文本路径差异，不作为全国性因果识别。***、**、* 分别表示在 1\%、5\%、10\% 水平显著。
    \end{flushleft}
\end{table}

政策文本证据的功能是展示高煤炭地区在政策约束后的地方承接表现，不用于识别绿色信贷的全国平均因果效应。山西和内蒙古均受到高煤炭体系约束，但二者在政策议程中的注意力配置不同：山西更偏向存量煤炭体系清洁化，政策重点围绕煤炭安全、先进产能、煤电改造、污染治理和煤炭产业链内部治理展开；内蒙古则在煤炭清洁治理之外更容易形成风光资源、外送通道、新能源基地、绿电和储能等新能源增量承接叙事。两种路径均属于绿色金融约束下的地方转型反应，但结构性治理能否发生，取决于地方资源禀赋、基础设施、项目体系和政策连续性。

\section{机制补充：政策文本、治理周期与承接稳定性}

\subsection{政策文本强度的构造与解释边界}

政策文本变量用于刻画地方政府在不同转型议题上的政策注意力和后验承接表现。本文对每篇政策文件进行简单关键词计数，不使用分词模型或 TF--IDF。具体而言，对每篇政策文件的标题、主题和正文进行拼接，然后对每个主题词表逐个统计出现次数。若一篇文件中 ``新能源'' 出现 5 次、``风电'' 出现 2 次、``光伏'' 出现 3 次、``储能'' 出现 1 次，则该文件的新能源主题词频为 11。年度层面再对所有政策文件加总，并构造每万字词频和政策覆盖比例。

这种构造方法的优点是透明、可复核，适合用于刻画政策注意力。其解释边界也需要明确：关键词计数不能识别语义方向，``限制煤炭'' 和 ``发展煤炭'' 都会增加煤炭相关词频。因此本文将其称为政策文本强度、政策注意力或后验承接表现，不将其解释为政策支持强度或因果机制强度。政策文本证据展示地方政府是否持续把绿色金融、煤炭清洁治理和新能源替代纳入政策议程，并不替代全国连续 DID 的识别。

\subsection{省级领导更替作为政策承接稳定性的补充解释}

全国绿色信贷政策需要通过地方政府政策议程、项目安排和产业治理体系才能落地。地方领导班子更替可能改变政策承接节奏、治理重点和产业政策表述，从而影响政策文本强度的阶段性变化。本文整理了山西和内蒙古省级党委、政府主要负责人更替节点，并将其与政策文本强度趋势进行对照。该材料作为附录补充证据，用于解释政策文本变化的阶段性，不作为独立因果识别。

山西在 2014 年、2016 年、2019--2021 年和 2022--2023 年出现多次党委或政府主要负责人更替。其中 2014 年王儒林任山西省委书记，与系统性腐败治理和煤炭治理背景高度相关；2016 年党政两端重组；2021 年以后又进入新的环保治理和能源安全叙事阶段。这些节点为山西煤炭清洁治理文本持续高位并阶段性重组提供了政治经济背景。

内蒙古的关键节点包括 2012 年党委书记更替、2016 年党政班子重组、2019 年党委书记更替、2021 年政府主要负责人变化和 2022 年党委书记更替。2012 年节点与绿色信贷政策出台高度接近，2016 年节点对应后续新能源项目和绿色金融政策承接期，2019 年节点又连接到涉煤腐败倒查和煤炭治理叙事。这些节点有助于解释内蒙古文本中绿色金融、煤炭清洁和新能源主题的阶段性抬升。

附录图将三类政策文本强度与领导更替年份同时展示。结果显示，若干文本强度跳升或政策叙事调整发生在地方领导班子更替或重组前后。这一证据支持一种有边界的机制解释：全国绿色信贷政策在地方层面的承接具有阶段性，可能受到地方治理周期、政策议程重组和产业转型任务再定位的共同影响。该机制不能替代主回归，也不能被解释为领导更替导致政策文本变化；它提供的是晋蒙路径差异背后的政策承接稳定性背景。

\section{结论与政策含义}

本文的证据链支持一个分层结论。绿色信贷政策确实改变了高煤炭地区，但这种改变首先体现为能源强度下降、煤炭使用强度下降和污染治理改善。煤炭消费占比也出现下降，说明约束性治理已经向能源使用结构产生外溢；但二氧化碳排放、火电装机占比和火电发电占比没有同步稳定调整，说明煤炭占比下降是约束性治理向结构性治理过渡的信号，尚不足以作为结构性治理完成的证据。

结构性治理需要更复杂的地方承接能力。低碳禀赋只是结构性治理的前置条件，尚不足以构成充分条件。真正的地方承接能力包括风光基础，也包括地方能否将金融约束转化为项目建设、并网消纳、外送通道和持续的低碳项目体系。山西和内蒙古展示了两种高煤炭地区转型路径。山西更偏向存量煤炭体系清洁化，通过煤炭安全、先进产能、煤电改造、污染治理和落后产能退出延续转型叙事。内蒙古在煤炭清洁治理之外，更容易形成风光基地、绿电外送、储能、源网荷储和新能源增量替代叙事，但这种替代仍需要时间、基础设施和政策连续性承接。

本文的政策意义在于，不同地区的绿色金融和转型金融工具应根据其所处的转型阶段进行匹配。对于存量煤炭锁定型地区，例如山西，更适合先推动煤炭体系内部清洁化、能效提升、污染治理、煤电改造、落后产能退出和有条件转型融资。对于煤炭基础仍重、但低碳资源和外送条件较好的地区，例如内蒙古，绿色信贷约束更可能与风光基地、绿电外送、储能、源网荷储和新能源增量替代结合。

绿色金融不能只按项目是否已经绿色筛选，否则会变成只服务绿色地区的金融工具。高煤炭地区真正需要的是转型金融组合：限制、改造、替代和退出并行。绿色信贷对高煤炭地区的意义，在于启动约束性治理并暴露结构性治理所需的地方承接条件；直接完成结构性低碳转型的证据不足。对于高煤炭地区而言，绿色金融政策的关键包括强化约束，也包括将约束与转型金融、低碳基础设施和地方政策承接结合起来，使治理性调整有可能进一步转化为结构性替代。

\end{document}
